\section{讨论与局限}
\label{sec:discussion}

\subsection{本复现的已知偏差}

如实记录与原论文设置的差异，是复现研究的基本要求。本次复现存在四处偏差，
按影响程度排序如下。

\paragraph{1. Stage 1 未按原方案使用 LLM 全量打标（影响最大）}
原论文用 LLM 为 5\,478 个正例事件对逐对标注支撑句，本文因平台免费额度耗尽，
改用以"事件所在句"为核心的规则法。第~\ref{sec:pseudo} 节的一致性对照表明
该替代方案在精确率上几乎无损（0.985）、在 97.4\% 的样本上被 LLM 答案完全包含，
但召回只有 0.753。这意味着 CGL 接收到的正样本偏少、负样本偏多，
其梯度方向比原方案更"保守"。

从结果看，这一偏差的后果集中在跨句流：句内流的伪标签与 LLM 高度一致（$P=0.994$），
对应 $F_1$ 提升 6.2 个点；跨句流的伪标签 $F_1$ 仅 0.746，对应 $F_1$ 与 iLIF 持平。
\textbf{换言之，本文实测到的"增益不均衡"很可能部分来自伪标签质量的不均衡，
而非方法本身的性质。}要分离这两个因素，需要在跨句流上补齐全量 LLM 标注后重跑，
这是最值得优先补做的实验。

\paragraph{2. 原文不可得，无法逐表核对}
CogECI 发表于 Elsevier 期刊，本次未能取得全文。因此本文的方法描述完全依据官方代码重建，
表中"iLIF 的 62.5"是公开可查的参照值，而非 CogECI 原文的主表数字。
若原文实际报告的是方向口径（三分类）结果，则本文的 0.6493 不可与之直接比较。
第~\ref{sec:protocol} 节用三条证据判定默认配置为存在性口径，但这是\emph{推断}，不是原文声明。

\paragraph{3. 单种子、单一划分}
全部实验固定 \texttt{seed=209}。5 折之间的标准差（$F_1$ 为 0.0261）反映的是折间差异，
不是种子间方差。要给出更可信的置信区间，至少需要 3--5 个种子重复。

\paragraph{4. 未做方向口径实验}
仓库同时提供了三分类的 pkl，本文只跑了二分类。方向口径的参照值是 iLIF 的 47.8，
若时间允许，这是一个成本低、信息量大的补充实验。

\subsection{关于 CGL 作用的观察}

消融实验（第~\ref{sec:ablation} 节）给出的是 CGL 权重取 0 / 3 / 6 的三点比较。
需要强调的是，CGL 的收益机制与伪标签质量强耦合：当伪标签在 97\% 以上的样本上
等于"两个事件所在句"时，CGL 实际上在教模型一件事——\emph{把事件对表示拉向它自身所在的两句}。
这在句内场景几乎等价于"事件对表示应当与承载它的句子语义一致"，是个很强的归纳偏置；
但在跨句场景，真正的证据往往分布在两句之外（如中间句交代了因果链），
此时"拉向两个事件句"就未必指向正确方向。

这一分析也解释了表~\ref{tab:reltype} 中一个乍看矛盾的现象：规则 C（区间 $\pm1$ 扩窗）
的召回高达 0.840 却使 $F_1$ 下降。扩窗把更多"正确答案"纳入伪标签，但这些额外句子
并非因果关系的承载句，CGL 会把事件对表示推向无关内容，反而损害表示质量。

\subsection{工程层面的反思}

本项目在工程上投入的时间明显多于算法设计，原因值得记录：

\begin{enumerate}
  \item \textbf{仓库不完整比算法难复现更常见。}缺失 \texttt{utils.py} 会直接报错，
        尚属"友好"；缺失伪标签 CSV 却静默回退全零，会在几十小时后给出一个"看起来正常"
        的低分结果，且很难定位。这类\emph{静默降级}是复现中最危险的一类问题，
        因此本文专门写了训练前终检脚本。
  \item \textbf{性能问题容易被误判为算力不足。}最初把 4.4 秒/文档归因于"6\,GB 卡太小"，
        但实测 GPU 利用率只有 8\%，说明瓶颈是调度而非算力。若不看利用率而直接下单租卡，
        成本会白花。
  \item \textbf{性能改写必须伴随等价性验证。}5 处改写里有 2 处（删除死代码、改选边实现）
        一旦出错不会报错、只会悄悄改变结果。逐项比对梯度与选边结果，是唯一可靠的验收方式。
\end{enumerate}

\subsection{合规性}

课程要求"检索库与提示不得包含测试标签或测试样本"。本文的实现满足该要求：
伪标签按文档独立生成；5 折中测试折文档不参与梯度更新；
所有 LLM 调用的提示词与原始响应均完整留痕（\texttt{raw\_*.jsonl}），可逐条审计。
需要如实说明的是，本文沿用了原论文的"label-informed"设定——提示词中告知 LLM
"该事件对确有因果"——这是原方法的设计，不是本项目的额外引入。
